\section{Introduction}
\IEEEPARstart{F}{uture} wireless networks are increasingly described not as fixed
cells but as graphs that change over time. A LEO mega-constellation is the clearest
example: hundreds to thousands of satellites move on known orbits, the inter
satellite links form and break on a schedule, and the routing, traffic, and resource
decisions all live on this moving graph. Unmanned aerial vehicle swarms, vehicular
ad-hoc networks, and software-reconfigurable fabrics share the same shape, a node
set with a time-varying edge set $G_t=(V,E_t)$. It is therefore natural that graph
neural networks (GNNs)~\cite{kipf2017gcn,wu2021gnnsurvey} have moved to the center of
learning-based network management, and equally natural that the community is asking
what quantum and quantum-inspired learning~\cite{biamonte2017qml} can add on top.

The special-issue theme of AI and quantum-enabled learning for future wireless
networks captures this momentum. The accompanying literature, however, leaves a
practitioner with two recurring gaps. First, the relationship between a classical
GNN and a quantum graph model is usually described by analogy rather than as a
precise statement, so it is hard to say what a quantum layer actually changes.
Second, quantum results are often reported in a way that is difficult to trust:
a single seed, a favorable operating point, or a metric whose scale is coupled to
the problem size. A reader is left unable to separate a genuine quantum effect from
an artifact of how the experiment was set up.

\textbf{This tutorial closes both gaps with one idea.} Classical message passing,
graph diffusion, and continuous-time quantum walks are not three unrelated methods;
they are three choices of a single \emph{propagation operator} $\Prop(\Lap)$ acting
on a graph signal through the Laplacian $\Lap$. Fixing the rest of the network and
varying only $\Prop$ makes the models param-matched and turns vague comparisons into
a controlled experiment. The operator view also makes the quantum mechanism concrete:
the heat kernel $e^{-t\Lap}$ spreads diffusively, with second moment growing linearly
in time, whereas the quantum walk $e^{-it\Lap}$ spreads \emph{ballistically}, with
second moment growing quadratically. That single difference is the whole of what
``going quantum'' buys at the propagation level, and we can measure it directly.

\textbf{Scope.} This is a tutorial and a focused survey, built around an original,
fully reproducible case study. The survey is organized by the propagation operator
rather than by author or year: we map classical, quantum, and quantum-inspired graph learning and its use in
wireless and non-terrestrial networks onto one axis (Section~\ref{sec:survey}), with the
depth deliberately uneven, favoring the operators we then build and evaluate. The
tutorial half teaches that one coherent framework, shows how to build a model from it,
and, crucially, how to evaluate it without deceiving ourselves. The case study is
deliberately small and fully reproducible so that every number and figure can be
regenerated from the companion repository. It is also independent of any specific
deployment pipeline: we study a clean node-potential task on a dynamic LEO graph, not a
particular routing or traffic-engineering system.

\textbf{Contributions.} The tutorial is organized around five takeaways.
\begin{itemize}
\item \textbf{C1, a unifying operator view and an operator-organized survey.} We present
GCN, heat diffusion, and the quantum walk as instances of $\Prop(\Lap)$, with diffusive
versus ballistic spreading as the distinguishing physics, and we survey the classical,
quantum, quantum-inspired, and dynamic-graph branches plus their wireless and
non-terrestrial applications on this single axis
(Sections~\ref{sec:survey},~\ref{sec:classical}--\ref{sec:quantum}).
\item \textbf{C2, build it.} We derive a param-matched quantum-inspired GNN layer that
runs on classical hardware, with pseudocode and reference code (Section~\ref{sec:buildit}).
\item \textbf{C3, a reproducible case study with scale and noise.} On a dynamic LEO graph
we compare the operators with multiple seeds and scale-free metrics, report the outcome
honestly, scale up to a $1584$-satellite shell, and add a NISQ-noise study showing how
device noise erodes the ballistic mechanism (Section~\ref{sec:casestudy}).
\item \textbf{C4, an honest-evaluation protocol.} We reproduce and then kill two
pitfalls that fabricate advantages, a scale-coupled normalization and single-seed
reporting, and give a reusable checklist (Section~\ref{sec:pitfalls}).
\item \textbf{C5, a decision guide and roadmap.} We say when to reach for each
operator and what stands between quantum-inspired models and true quantum hardware
(Sections~\ref{sec:decision}--\ref{sec:roadmap}).
\end{itemize}

\textbf{A note on honesty.} Our case study does not hand the quantum-inspired operator
an easy win, and it shows why the evaluation discipline matters. At small scale the
quantum-inspired walk loses to a classical global operator and is noisier; at medium
scale it only ties. Its advantage appears only at the largest scale we test ($1584$
satellites), where it becomes the best operator, the diameter-driven trend its ballistic
spreading predicts, and even there we temper the claim with a three-seed caveat. We
report all of this plainly. For a tutorial, teaching readers how to reach and recognize a
scale-dependent, conditional advantage, and how a single seed or scale would have
misled, is at least as valuable as a headline win.

Fig.~\ref{fig:reading} shows how to read the paper depending on your goal.
\begin{figure}[t]
\centering
\resizebox{\columnwidth}{!}{%
\begin{tikzpicture}[node distance=4mm]
\node[box] (goal) {Your goal};
\node[hot, right=8mm of goal] (mech) {Understand the\\mechanism};
\node[hot, below=3mm of mech] (build) {Build a model};
\node[hot, below=3mm of build] (eval) {Evaluate it\\honestly};
\node[plain, right=8mm of mech] (s34) {Sec.~\ref{sec:classical}--\ref{sec:quantum}};
\node[plain, right=8mm of build] (s5) {Sec.~\ref{sec:buildit}--\ref{sec:casestudy}};
\node[plain, right=8mm of eval] (s7) {Sec.~\ref{sec:pitfalls}};
\draw[flow] (goal.east) -- (mech.west);
\draw[flow] (goal.east) |- (build.west);
\draw[flow] (goal.east) |- (eval.west);
\draw[flow] (mech) -- (s34);
\draw[flow] (build) -- (s5);
\draw[flow] (eval) -- (s7);
\end{tikzpicture}}
\caption{Reading path. The unifying mechanism is in
Sections~\ref{sec:classical}--\ref{sec:quantum}; how to build a model in
Sections~\ref{sec:buildit}--\ref{sec:casestudy}; how to evaluate it without fooling
yourself in Section~\ref{sec:pitfalls}.}
\label{fig:reading}
\end{figure}

